% GENERATED FILE. Edit canonical lessons, then run npm run generate:books.
% canonical-source-hash: fnv1a64:88a72add3e73e59a
% canonical-lessons: AR-C79-baqara, AR-C79-maiz, AR-C79-kharuf, AR-C79-far, AR-C79-asad

\chapter{Cow, Goat, Sheep, Mouse, Lion}
\label{ch:ar-cow-goat-sheep-mouse-lion}
\hlchaptermodality{79}

\begin{chapteropening}
\textbf{By the end of this chapter:} \emph{I can say a cow, a goat, a sheep, a mouse and a lion.}

Complete the last lesson of chapter 79: I can say a cow, a goat, a sheep, a mouse and a lion.
\end{chapteropening}

\section[\texorpdfstring{\ar{بقرة} (baqara)}{بقرة (baqara)}]{\ar{بقرة} (baqara) — a cow}

\label{lesson:AR-C79-baqara}



\begin{quote}
Before the new one: say the Arabic for the sky, then the Arabic for a horse.
\end{quote}

\subsection*{You'll want to know: \ar{بقرة}}
\textbf{\ar{بقرة}} — \emph{baqara} — \textquotedblleft{}a cow\textquotedblright{}.

\begin{cousinweb}
\textbf{\ar{بقرة}} (\emph{baqara}) is built on the native Arabic root \emph{b-q-r}, \textquotedblleft{}to slit, to split\textquotedblright{} (Lane). \emph{Baqara} is one head of \emph{baqar}, \textquotedblleft{}cattle\textquotedblright{}, which Lane lists under this root.
\end{cousinweb}

\begin{grammarlens}[title={What you've built}]
One more word for this chapter.
\end{grammarlens}

\subsection*{Your turn}
\begin{itemize}
\raggedright
  \item \emph{Say it:} \emph{baqara}
  \item \emph{Say it:} \emph{baqara}, once more
  \item \emph{Say it:} say \emph{baqara}
  \item \emph{Recall:} say the Arabic for the sky, then the Arabic for a horse, then say \emph{baqara} again
\end{itemize}

\begin{tcolorbox}[breakable,colback=teal!4,colframe=teal!35!black,title={Before you move on}]
What does \ar{بقرة} mean? (\textquotedblleft{}A cow\textquotedblright{}.) Say it once more.
\end{tcolorbox}
\section[\texorpdfstring{\ar{ماعز} (māʿiz)}{ماعز (māʿiz)}]{\ar{ماعز} (māʿiz) — a goat}

\label{lesson:AR-C79-maiz}



\begin{quote}
Before the new one: say the Arabic for a horse, then the Arabic for a cow.
\end{quote}

\subsection*{You'll want to know: \ar{ماعز}}
\textbf{\ar{ماعز}} — \emph{māʿiz} — \textquotedblleft{}a goat\textquotedblright{}.

\begin{cousinweb}
\textbf{\ar{ماعز}} (\emph{māʿiz}) is built on the native Arabic root \emph{m-ʿ-z}, \textquotedblleft{}goats\textquotedblright{} (Lane). \emph{Māʿiz} is one goat; \emph{maʿz} is the collective.
\end{cousinweb}

\begin{grammarlens}[title={What you've built}]
One more word for this chapter.
\end{grammarlens}

\subsection*{Your turn}
\begin{itemize}
\raggedright
  \item \emph{Say it:} \emph{māʿiz}
  \item \emph{Say it:} \emph{māʿiz}, once more
  \item \emph{Say it:} say \emph{māʿiz}
  \item \emph{Recall:} say the Arabic for a horse, then the Arabic for a cow, then say \emph{māʿiz} again
\end{itemize}

\begin{tcolorbox}[breakable,colback=teal!4,colframe=teal!35!black,title={Before you move on}]
What does \ar{ماعز} mean? (\textquotedblleft{}A goat\textquotedblright{}.) Say it once more.
\end{tcolorbox}
\section[\texorpdfstring{\ar{خروف} (kharūf)}{خروف (kharūf)}]{\ar{خروف} (kharūf) — a sheep}

\label{lesson:AR-C79-kharuf}



\begin{quote}
Before the new one: say the Arabic for a cow, then the Arabic for a goat.
\end{quote}

\subsection*{You'll want to know: \ar{خروف}}
\textbf{\ar{خروف}} — \emph{kharūf} — \textquotedblleft{}a sheep\textquotedblright{}.

\begin{cousinweb}
\textbf{\ar{خروف}} (\emph{kharūf}) is built on the native Arabic root \emph{kh-r-f}, \textquotedblleft{}to gather fruit; to graze\textquotedblright{} (Lane). Lane's \emph{kharūf} is a lamb, said to be so called because it grazes here and there; the same root gives \emph{kharīf}, \textquotedblleft{}autumn\textquotedblright{}.
\end{cousinweb}

\begin{grammarlens}[title={What you've built}]
One more word for this chapter.
\end{grammarlens}

\subsection*{Your turn}
\begin{itemize}
\raggedright
  \item \emph{Say it:} \emph{kharūf}
  \item \emph{Say it:} \emph{kharūf}, once more
  \item \emph{Say it:} say \emph{kharūf}
  \item \emph{Recall:} say the Arabic for a cow, then the Arabic for a goat, then say \emph{kharūf} again
\end{itemize}

\begin{tcolorbox}[breakable,colback=teal!4,colframe=teal!35!black,title={Before you move on}]
What does \ar{خروف} mean? (\textquotedblleft{}A sheep\textquotedblright{}.) Say it once more.
\end{tcolorbox}
\section[\texorpdfstring{\ar{فأر} (faʾr)}{فأر (faʾr)}]{\ar{فأر} (faʾr) — a mouse}

\label{lesson:AR-C79-far}



\begin{quote}
Before the new one: say the Arabic for a goat, then the Arabic for a sheep.
\end{quote}

\subsection*{You'll want to know: \ar{فأر}}
\textbf{\ar{فأر}} — \emph{faʾr} — \textquotedblleft{}a mouse\textquotedblright{}.

\begin{cousinweb}
\textbf{\ar{فأر}} (\emph{faʾr}) is built on the native Arabic root \emph{f-ʾ-r}, \textquotedblleft{}rats, mice\textquotedblright{} (Lane). An old noun with its own root in Lane.
\end{cousinweb}

\begin{grammarlens}[title={What you've built}]
One more word for this chapter.
\end{grammarlens}

\subsection*{Your turn}
\begin{itemize}
\raggedright
  \item \emph{Say it:} \emph{faʾr}
  \item \emph{Say it:} \emph{faʾr}, once more
  \item \emph{Say it:} say \emph{faʾr}
  \item \emph{Recall:} say the Arabic for a goat, then the Arabic for a sheep, then say \emph{faʾr} again
\end{itemize}

\begin{tcolorbox}[breakable,colback=teal!4,colframe=teal!35!black,title={Before you move on}]
What does \ar{فأر} mean? (\textquotedblleft{}A mouse\textquotedblright{}.) Say it once more.
\end{tcolorbox}
\section[\texorpdfstring{\ar{أسد} (asad)}{أسد (asad)}]{\ar{أسد} (asad) — a lion}

\label{lesson:AR-C79-asad}



\begin{quote}
Before the new one: say the Arabic for a sheep, then the Arabic for a mouse.
\end{quote}

\subsection*{You'll want to know: \ar{أسد}}
\textbf{\ar{أسد}} — \emph{asad} — \textquotedblleft{}a lion\textquotedblright{}.

\begin{cousinweb}
\textbf{\ar{أسد}} (\emph{asad}) is built on the native Arabic root \emph{ʾ-s-d}, \textquotedblleft{}a lion\textquotedblright{} (Lane). An old noun with its own root in Lane.
\end{cousinweb}

\begin{grammarlens}[title={What you've built}]
One more word for this chapter.
\end{grammarlens}

\subsection*{Your turn}
\begin{itemize}
\raggedright
  \item \emph{Say it:} \emph{asad}
  \item \emph{Say it:} \emph{asad}, once more
  \item \emph{Say it:} say \emph{asad}
  \item \emph{Recall:} say the Arabic for a sheep, then the Arabic for a mouse, then say \emph{asad} again
\end{itemize}

\begin{tcolorbox}[breakable,colback=teal!4,colframe=teal!35!black,title={Before you move on}]
What does \ar{أسد} mean? (\textquotedblleft{}A lion\textquotedblright{}.) Say it once more.
\end{tcolorbox}
